\section*{Capítulo \ref{ch:b3:genetic-engineering} --- Ingeniería genética y biotecnología}

\begin{solution}{exo:b3:genetic-engineering:1}
Un origen de replicación, para que el hospedador copie el \omterm{def:b3:genetic-engineering:tools}{plásmido} y lo
pase a las células hijas; un \omterm{def:b3:genetic-engineering:tools}{marcador de selección}, por lo común una
resistencia a un antibiótico, para que solo crezcan las células con
\omterm{def:b3:genetic-engineering:tools}{plásmido}; y un sitio de restricción único (o un grupo de sitios) en el
que se liga el inserto sin cortar el \omterm{def:b3:genetic-engineering:tools}{plásmido} por otro lado.
\end{solution}

\begin{solution}{exo:b3:genetic-engineering:2}
Una secuencia dada de seis bases aparece una vez cada $4^{6} = 4096$ bp
en ADN aleatorio: unos $4.6\times 10^{6}/4096 \approx 1100$ fragmentos de
\qty{4.1}{kb} de media. Los genomas reales no son aleatorios --- la
composición de bases, el uso de codones y la evitación de algunas
secuencias desplazan el recuento (el número real de sitios EcoRI de
\emph{E.~coli} es de varios centenares).
\end{solution}

\begin{solution}{exo:b3:genetic-engineering:3}
La desnaturalización a \qty{95}{\celsius} separa las hebras; la
hibridación a \qtyrange{50}{65}{\celsius} deja que los \omterm{def:b3:genetic-engineering:pcr}{cebadores} se
apareen con sus sitios; y la extensión a \qty{72}{\celsius} deja que la
polimerasa copie desde cada \omterm{def:b3:genetic-engineering:pcr}{cebador}. Una polimerasa corriente se
destruiría a \qty{95}{\celsius} en el primer ciclo y habría que añadirla
de nuevo treinta veces; la Taq sobrevive a todos los ciclos.
\end{solution}

\begin{solution}{exo:b3:genetic-engineering:4}
Un \omterm{def:b3:genetic-engineering:crispr}{ARN guía} cuyas veinte bases correspondan a la diana, y un \omterm{def:b3:genetic-engineering:crispr}{PAM} (NGG)
inmediatamente en 3$'$ de la diana en la hebra no diana. Tras la \omterm{def:b3:dna-repair:lesions}{rotura
de doble hebra}: la unión de extremos, que deja una inserción o deleción
pequeña y suele destruir el marco de lectura del gen; o la \omterm{def:b3:dna-repair:dsb}{recombinación
homóloga} con un molde que lleve la secuencia deseada, que la escribe.
\end{solution}

\begin{solution}{exo:b3:genetic-engineering:5}
$50\times 1.9^{35} = 50\times 5.7\times 10^{9} \approx 3\times 10^{11}$
copias. $\Delta C_{T} = 7.3$; razón $1.95^{7.3} = e^{7.3\ln 1.95}
\approx 130$: la primera muestra tenía unas $130$ veces más molde.
\end{solution}

\begin{solution}{exo:b3:genetic-engineering:6}
Las bacterias no pueden cortar y empalmar: una copia genómica con
intrones se transcribiría en un mensaje inservible, de modo que se usa el
ADNc sin intrones, copiado del ARNm maduro. Otros obstáculos: la
polimerasa bacteriana no reconoce el promotor humano (aportar un promotor
bacteriano y un sitio de unión al ribosoma); el uso de codones humano
difiere (sintetizar el gen con los codones preferidos del hospedador); la
proteína puede no plegarse o agregarse (bajar la temperatura, fusionarla
a un compañero soluble o replegarla desde los cuerpos de inclusión);
falta la glicosilación (usar levadura o células de mamífero si importa);
y los puentes disulfuro necesitan el periplasma oxidante.
\end{solution}

\begin{solution}{exo:b3:genetic-engineering:7}
Las células madre dirigidas se inyectan en un blastocisto hospedador y se
mezclan con sus propias células, de modo que el ratón está construido a
partir de dos genotipos --- una quimera --- y solo transmite el alelo si
algunas de sus células germinales derivan de las células dirigidas.
Cruzar la quimera con un ratón silvestre da heterocigotos (a partir de la
mitad de sus gametos si la línea germinal está dirigida a medias);
intercruzar heterocigotos da una cuarta parte de homocigotos para la
supresión.
\end{solution}

\begin{solution}{exo:b3:genetic-engineering:8}
Exactas: $6.4\times 10^{9}\times 4^{-20}\times (1/16) \approx 4\times
10^{-4}$ sitios (el GG del \omterm{def:b3:genetic-engineering:crispr}{PAM} cuesta un factor $16$) --- ninguno, en la
práctica. A dos discrepancias hay $1 + 60 + 1710 = 1771$ secuencias:
$1771\times 4\times 10^{-4} \approx 0.7$ sitios por azar. El genoma real
no es aleatorio, y una guía se contrasta directamente con él.
\end{solution}

\begin{solution}{exo:b3:genetic-engineering:9}
$p_{1} = 0.05 + 0.8\times 0.05\times 0.95 = 0.088$; $p_{2} = 0.152$;
$p_{3} = 0.255$; luego $0.41$ y $0.60$: unas cinco generaciones para
pasar de $0.5$, frente a la estimación
$\ln(1/(2\times 0.05))/\ln 1.8 \approx 4$.
\end{solution}

\begin{solution}{exo:b3:genetic-engineering:10}
La unión de extremos actúa en toda célula, incluidas las células madre
quiescentes, en las que la \omterm{def:b3:dna-repair:dsb}{recombinación homóloga}, un proceso de S y G2,
es poco eficiente; no exige entregar ningún molde y su producto ---
cualquier interrupción del potenciador --- hace el trabajo, mientras que
la corrección exige la secuencia exacta y rinde una minoría de células.
Además, la misma edición sirve para toda mutación de la
$\beta$-globina, falciforme o talasémica. Inconveniente: elimina un
elemento regulador (con todas las demás funciones que tenga), deja el
alelo falciforme donde estaba y la mezcla de inserciones y deleciones no
está controlada.
\end{solution}

\begin{solution}{exo:b3:genetic-engineering:11}
Con una eficacia biológica de los homocigotos $1 - s$,
$p' = \bigl[p^{2}(1-s) + p(1-p)(1+c)
\bigr]/(1 - sp^{2})$. Desde la rareza los homocigotos son despreciables y
el impulsor crece como $(1+c)p$ sea cual sea $s$; que llegue a la
fijación se decide cerca de $p = 1$, donde un alelo silvestre raro de
frecuencia $q$ vuelve como $q' \approx q(1-c)/(1-s)$: hay fijación si
$c > s$, y un equilibrio intermedio en caso contrario. Resistencia: cada
conversión fallida (unión de extremos en lugar de recombinación) deja un
alelo que la guía ya no reconoce; si esos alelos son biológicamente
eficaces --- una inserción o deleción en fase en una región no esencial
---, no tienen coste mientras que el impulsor sí lo tiene, de modo que la
selección los favorece y sustituyen al impulsor. Los remedios son dirigir
el corte a secuencias esenciales donde las inserciones y deleciones sean
letales, y usar varias guías a la vez.
\end{solution}

\begin{solution}{exo:b3:genetic-engineering:12}
La \omterm{def:b3:genetic-engineering:crispr}{Cas9} inyectada en un cigoto puede actuar después de las primeras
divisiones, de modo que células distintas reciben reparaciones distintas
--- el mosaicismo ---; la reparación de cada rotura la elige la célula, y
la unión de extremos da inserciones y deleciones impredecibles, grandes
deleciones o pérdida de heterocigosidad en el corte; el resultado por
\omterm{def:b3:dna-repair:dsb}{recombinación homóloga} es minoritario; hay roturas \omterm{prop:b3:genetic-engineering:editors}{fuera de diana} en
sitios que no pueden predecirse todos; y el embrión solo puede
comprobarse secuenciando unas pocas células de biopsia, que no pueden
hablar por las demás. Pruebas necesarias: métodos que editen todas las
células por igual (antes de la primera fase S) con eficiencias cercanas
al \qty{100}{\%}, la secuenciación del genoma completo de cada célula de
embriones animales editados que no muestre ningún cambio no pretendido, y
el seguimiento a largo plazo de animales editados a lo largo de
generaciones --- nada de lo cual existe.
\end{solution}

\begin{solution}{pb:b3:genetic-engineering:1}
\textbf{1.} Una vez cada $4096$ bp. $P(\text{sin sitio}) = (1 - 1/4096)^{1500}
= e^{-0.37} = 0.69$.
\textbf{2.} Elíjase una enzima sin sitio en el gen; o amplifíquese el gen
por PCR con \omterm{def:b3:genetic-engineering:pcr}{cebadores} que lleven sitios de restricción nuevos en sus
extremos 5$'$ (o úsese clonación de extremos romos o independiente de
ligación).
\textbf{3.} \qty{0.1}{\micro g}\,$\times 2\times 10^{7} = 2\times
10^{6}$ colonias; el \qty{10}{\%}, $2\times 10^{5}$, llevan el inserto.
\textbf{4.} Blancas: el \qty{10}{\%}. La mitad de los insertos están en
la orientación correcta, de modo que
$P(\text{ninguna en } n \text{ blancas}) = 0.5^{n} \le 0.01$ da
$n \ge 6.6$: tómense siete.
\textbf{5.} $32$ duplicaciones, $2^{32} = 4.3\times 10^{9}$ células,
$2.1\times
10^{11}$ \omterm{def:b3:genetic-engineering:tools}{plásmidos}; cada uno $5500\times 650 = 3.6\times 10^{6}$ Da $=
5.9\times 10^{-18}$ g: \qty{1.3}{\micro g} de \omterm{def:b3:genetic-engineering:tools}{plásmido}.
\textbf{6.} La clonación solo necesita replicación, que aporta el origen
del \omterm{def:b3:genetic-engineering:tools}{plásmido}; la ARN-polimerasa bacteriana no lee el promotor humano, de
modo que para la expresión hay que poner delante de la secuencia
codificante (sin intrones) un promotor bacteriano y un sitio de unión al
ribosoma.
\textbf{7.} $N_{0} = N_{T}/2^{C_{T}}$: $10^{12}/2^{28} \approx 3700$
copias; con $C_{T} = 35$, $10^{12}/2^{35} \approx 29$ copias.
\textbf{8.} $2^{n} = 10^{12}$: $n = 39.9$, cuarenta ciclos.
\textbf{9.} $10^{12}/2^{38} \approx 3.6$ copias. Más allá de unos $40$
ciclos, una sola molécula contaminante, o los artefactos de los
\omterm{def:b3:genetic-engineering:pcr}{cebadores}, llegan al umbral, y una muestra de menos de una copia no
significa nada.
\textbf{10.} Una molécula cruza en el ciclo $40$: un positivo. Se evita
con salas y material separados para montar la reacción y para manipular
los productos, un flujo de trabajo en un solo sentido, controles
negativos en cada carrera y la destrucción enzimática de los amplicones
arrastrados.
\textbf{11.} $3700/0.4 \approx 9000$ copias de ARN.
\textbf{12.} Informa $2^{3.3} \approx 10$ veces; la verdadera es
$1.9^{3.3} =
e^{3.3\ln 1.9} \approx 8.3$ veces.
\textbf{13.} $40\times \qty{35}{\micro g} = \qty{1.4}{mg}$ al día,
\qty{0.51}{g} al año; $\times 10^{7}$: \qty{5100}{kg} al año.
\textbf{14.} $5.1\times 10^{6}/5800 \approx 880$ mol; $5.3\times 10^{26}$
moléculas.
\textbf{15.} $5.1\times 10^{6}$ g $/ (20$ g/L$) = 2.6\times 10^{5}$ L
$= \qty{260}{m^3}$ --- una décima parte de la piscina.
\textbf{16.} $5100\times 8000 = 4.1\times 10^{7}$ kg de páncreas, a
\qty{0.1}{kg} cada uno: $4\times 10^{8}$ cerdos al año.
\textbf{17.} La hormona recombinante es idéntica a la humana, de modo que
no provoca anticuerpos ni reacciones alérgicas; no lleva patógenos
animales (virus, priones); el suministro no depende del sacrificio; y la
secuencia puede mejorarse.
\textbf{18.} Los residuos que tocan el receptor no cambian, de modo que
la unión y la señalización son las de la insulina; los residuos alterados
están en la superficie por la que las moléculas de insulina se asocian en
dímeros y hexámeros, y cambiarlos cambia la velocidad con la que el
depósito inyectado se deshace en monómeros absorbibles.
\textbf{19.} $6.4\times 10^{9}\times 4^{-20}/16 \approx 3.6\times
10^{-4}$: ninguna coincidencia exacta por azar.
\textbf{20.} $1 + 60 + 1710 + 27\times 1140 = \num{32551}$ secuencias;
$\num{32551}\times 3.6\times 10^{-4} \approx 12$ sitios \omterm{prop:b3:genetic-engineering:editors}{fuera de diana}
por azar.
\textbf{21.} $10^{7}\times 10^{-4} = 1000$ células. Una célula madre vive
y se divide durante toda la vida del paciente; una que haya perdido un
supresor de tumores, o adquirido una translocación, puede fundar un clon
que acabe en leucemia.
\textbf{22.} $p_{1} = 0.0019$, $p_{2} = 0.0036$, $p_{3} = 0.0069$; el
crecimiento geométrico por $1.9$ por generación da $\ln 500/\ln 1.9
\approx 10$; iterando la recurrencia, la frecuencia pasa de $0.5$ en la
generación $11$ ($0.45 \to 0.68$).
\textbf{23.} Alrededor de un año. Un coste del \qty{20}{\%} en los
homocigotos es menor que $c = 0.9$, de modo que el impulsor llega igual a
la fijación, más despacio, mientras la eficacia biológica media de la
población cae --- que es la finalidad de un impulsor de supresión; un
coste superior a $c$ lo detendría en una frecuencia intermedia.
\textbf{24.} $10^{6}\times 0.1 = 10^{5}$ alelos de resistencia en una
generación. Al no poder cortarse y ser, si son eficaces, gratuitos
mientras el impulsor es costoso, quedan favorecidos y se propagan, y el
impulsor se elimina --- de ahí que los impulsores se dirijan a sitios
donde los productos de la unión de extremos sean letales, y con varias
guías a la vez.
\textbf{25.} Unas $3700$ copias en la primera muestra; unos
\qty{260}{m^3} de fermentador al año para la insulina del mundo; y unas
$11$ generaciones, un año de mosquito, para el impulsor.
\end{solution}
